\documentclass[10pt]{article}
\usepackage[margin=2cm]{geometry}
\usepackage{booktabs}
\usepackage{graphicx}
\usepackage{hyperref}
\usepackage{amsmath}
\setlength{\emergencystretch}{3em}

\begin{document}

\title{CSCI 5942: AI Engineering Assignment 6\\Serving GPT-2}
\author{Bill Hsu}
\maketitle

% =====================================================================
\section*{Setup}

All runs used one NVIDIA A40 (46\,GB) on the NCSA Delta node
\texttt{gpub052} (partition \texttt{gpuA40x4}), with the
\texttt{pytorch-conda/2.12} module (PyTorch 2.12.1, CUDA 13.0). We ran
every command in one Slurm job, \texttt{slurm/run\_all.sbatch} (job
23390832). \texttt{sample.py} runs gpt2-xl in fp32; \texttt{bench\_kv.py}
and \texttt{bench\_serve.py} run it in bf16. Times are wall-clock with
\texttt{torch.cuda.synchronize()} before and after. In
\texttt{bench\_serve.py}, every request arrives at time 0, so latency is
the time from the start of the run until a request finishes, and mean
batch is the mean decode batch size. As in \texttt{bench\_serve.py},
1\,GB means $10^9$ bytes. All tests pass on the GPU node:
\texttt{tests/test\_model.py} (4 passed), \texttt{tests/test\_kvcache.py}
(8 passed), \texttt{tests/test\_paged.py} (8 passed) and
\texttt{tests/test\_engine.py} (7 passed).

% =====================================================================
\section*{Part 1: The KV Cache}

\begin{table}[htbp]
\caption{\texttt{sample.py} with gpt2-xl (fp32, batch 1, greedy decoding,
prompt \texttt{"The capital of France is"}). In both runs, the cached and
uncached tokens are the same.}
\label{tab:sample}
\centering
\begin{tabular}{rrrc}
\toprule
New tokens & Cached (ms/token) & Uncached (ms/token) & Tokens match \\
\midrule
200 & 22.2 & 33.2 & True \\
800 & 22.9 & 93.7 & True \\
\bottomrule
\end{tabular}
\end{table}

\paragraph{Q1: Why does the uncached time grow, and the cached time not?}
Without a cache, the model runs the whole sequence again for every new
token, so one step costs more as the sequence gets longer, and the total
cost grows as $O(N^2)$. Going from 200 to 800 new tokens, the uncached
time per token grows $2.8\times$, from 33.2 to 93.7\,ms. With a cache,
each step runs only the one new token. Attention still reads every cached
key and value, but that is small compared with reading the model weights,
so the cached time stays almost flat (22.2 to 22.9\,ms).
At 200 tokens, the gap is small (33.2 vs.\ 22.2\,ms) because decoding is
memory-bound: most of each step is spent reading the weights from GPU
memory, and that cost is the same for one token or for 200.

\paragraph{Q2: Size of the KV cache for one 1024-token request.}
For every token, each of the 48 layers stores one key and one value. Each
has $n_\text{embd} = n_\text{head} \times \text{head\_dim} = 25 \times 64
= 1600$ numbers, and each number is 2\,B in bf16:
\begin{align*}
\text{per token} &= 2\ (\text{K, V}) \times 48 \times 1600 \times 2\,\text{B}
  = 307{,}200\,\text{B} = 300\,\text{KiB}, \\
\text{per request} &= 307{,}200\,\text{B} \times 1024
  = 314{,}572{,}800\,\text{B} \approx 314.6\,\text{MB}.
\end{align*}
This is 300\,MiB, and the per-token size matches the
\texttt{300 KiB/token} that \texttt{bench\_kv.py} prints.

% =====================================================================
\section*{Part 2: PagedAttention}

\begin{table}[htbp]
\caption{\texttt{bench\_kv.py}: gpt2-xl in bf16, 32 prompts of 95--485
tokens (mean 267), 128 new tokens each. Used \% is the fraction of
reserved KV slots that hold a token. The first contiguous row reserves
the longest prompt plus 128 slots for every row; the second shows what
reserving 1024 slots per row would use.}
\label{tab:kv}
\centering
\begin{tabular}{lrrrrr}
\toprule
Cache & Prefill (ms) & Decode (ms/step) & Reserved & Used & Used \% \\
\midrule
contiguous                 & 609 & 22.9 & 19{,}616 & 12{,}616 & 64.3\% \\
contiguous (max\_len 1024) &     &      & 32{,}768 & 12{,}616 & 38.5\% \\
paged, block 8   & 855 & 36.5 & 12{,}728 & 12{,}616 & 99.1\% \\
paged, block 16  & 764 & 36.6 & 12{,}880 & 12{,}616 & 98.0\% \\
paged, block 32  & 770 & 37.0 & 13{,}120 & 12{,}616 & 96.2\% \\
paged, block 64  & 768 & 37.7 & 13{,}568 & 12{,}616 & 93.0\% \\
paged, block 128 & 769 & 39.4 & 14{,}336 & 12{,}616 & 88.0\% \\
paged, block 256 & 769 & 43.9 & 16{,}640 & 12{,}616 & 75.8\% \\
\bottomrule
\end{tabular}
\end{table}

Table~\ref{tab:kv} shows that paging wastes very little memory: at block
16, 98.0\% of the reserved slots hold a token, against 38.5\% for a
contiguous cache that reserves 1024 slots per request. The only waste is
the unused slots in each sequence's last block. The cost is speed: every
paged decode step is slower than the contiguous one (36.6 vs.\ 22.9\,ms at
block 16), because in every layer the gather copies each sequence's
blocks out of the pool into one contiguous tensor. A real system uses an
attention kernel that reads the blocks directly from the pool and does
not pay for this copy. Paged prefill is also slower (764 vs.\ 609\,ms),
mostly because \texttt{bench\_kv.py} prefills each paged prompt on its
own, while the contiguous cache prefills all 32 prompts in one padded
batch.

\paragraph{Q1: How many of these requests fit in 4\,GB?}
\begin{itemize}
\item \textbf{Contiguous.} Each request reserves 1024 slots, which is
  314.6\,MB (Part 1).
  $4 \times 10^9 / 314{,}572{,}800 = 12.7$, so only \textbf{12} requests
  fit.
\item \textbf{Paged, block 16.} 4\,GB holds
  $4 \times 10^9 / (16 \times 307{,}200) = 813$ blocks. At their longest,
  the 32 requests hold 12{,}880 slots (Table~\ref{tab:kv}), which is
  805 blocks, or about 25 blocks (123.6\,MB) per request. Since
  $805 \le 813$, \textbf{all 32} requests fit at the same time.
\end{itemize}
Based on our calculations, paging fits about $2.7\times$ more of these
requests in the same memory.

\paragraph{Q2: Which block size?}
I would use \textbf{block size 16}. It keeps 98.0\% of the reserved slots
in use, and its decode time is the same as block 8 (36.6 vs.\ 36.5\,ms).
Very small blocks make the block tables longer and add more indexing work
for every token. Very large blocks waste the empty slots of each
sequence's last block (only 75.8\% used at block 256), and decode is
slower too (43.9\,ms), because the gather also copies those empty slots.

% =====================================================================
\section*{Part 3: Continuous Batching}

\begin{table}[htbp]
\caption{\texttt{bench\_serve.py}: gpt2-xl in bf16, 400 requests arriving
at once, prompts of 16--64 tokens, outputs of 8--128 tokens (27{,}418 in
total), block size 16, max batch 64. Latency is in seconds.}
\label{tab:serve}
\centering
\small
\begin{tabular}{llrrrrrr}
\toprule
KV cache & Engine & Time (s) & Tokens/s & Mean lat. & p90 lat. & Mean batch & Preempt \\
\midrule
4\,GB (813 blocks) & static     & 29.6 &  926 & 14.6 & 25.4 & 31.1 & 0 \\
                   & continuous & 21.9 & 1250 & 12.0 & 19.8 & 54.0 & 0 \\
\midrule
1\,GB (203 blocks) & static     & 53.1 &  517 & 26.3 & 47.6 & 14.6 & 0 \\
                   & continuous & 37.8 &  726 & 19.1 & 33.9 & 22.5 & 0 \\
                   & continuous (optimistic) & 33.3 & 824 & 17.2 & 30.2 & 32.5 & 176 \\
\bottomrule
\end{tabular}
\end{table}

\paragraph{Q1: Static vs.\ continuous.}
Table~\ref{tab:serve} shows that continuous batching improves all three
metrics at both budgets.
\begin{itemize}
\item \textbf{Tokens/s} (926 to 1250, +35\%, at 4\,GB; 517 to 726, +40\%,
  at 1\,GB). Decoding is memory-bound, so a step costs about the same for
  a larger batch, and a larger batch makes more tokens per step.
\item \textbf{Mean latency} (14.6 to 12.0\,s, $-18$\%, at 4\,GB; 26.3 to
  19.1\,s, $-27$\%, at 1\,GB). A waiting request starts as soon as a slot
  is free, so it does not wait for the longest request of a whole static
  batch to finish.
\item \textbf{Mean batch} (31.1 to 54.0 at 4\,GB; 14.6 to 22.5 at 1\,GB).
  Finished requests are retired and new ones are admitted at every
  iteration, so their slots and blocks are reused right away instead of
  staying idle.
\end{itemize}

\paragraph{Q2: What did paging buy at 1\,GB?}
\begin{enumerate}
\item A contiguous cache must reserve 1024 slots, or 314.6\,MB, for every
  request (Part 1). $10^9 / 314{,}572{,}800 = 3.2$, so only \textbf{3}
  requests fit at once.
\item In Part 1 and Part 2, one gpt2-xl decode step with a contiguous
  cache takes about 22--23\,ms, and this stays almost the same from
  batch 1 (22.2\,ms) to batch 32 (22.9\,ms). With 3 requests, the
  throughput would be at most about
  $3 / 0.023\,\text{s} \approx$ \textbf{130 tokens/s}, before counting
  prefill. This is an estimate, not a measurement.
\item Our continuous run reaches 726 tokens/s at 1\,GB, about
  $5.6\times$ the estimate. Each paged step is slower because of the
  gather, but these requests are short (at most $64 + 128 = 192$ tokens,
  or 12 blocks), and paging gives them only the blocks they use. This
  lets about 22 requests run at once instead of 3, and the larger batch
  is worth much more than the cost of the gather.
\end{enumerate}

\paragraph{Q3: Optimistic vs.\ safe admission at 1\,GB.}
With optimistic admission, tokens/s goes from 726 to 824 (+13\%), and
the mean batch goes from 22.5 to 32.5 (+44\%). Mean latency also falls,
from 19.1 to 17.2\,s.
\begin{itemize}
\item \textbf{Benefit.} Safe admission reserves blocks for every
  request's full length, and most of those blocks stay empty for a long
  time. Optimistic admission reserves only the next token, so more
  requests run at once and the pool is used better. Also, recomputing is 
  cheaper than it looks: the engine re-runs the prompt and all generated
  tokens in one prefill, which is far more efficient than the decode steps
  that first produced them.
\item \textbf{Cost.} When the pool runs out, a request is preempted and
  must be recomputed from the start. There were 176 preemptions, and the
  prefill and decode work done for them before each preemption is
  wasted. A preempted request also waits again, so its latency is harder
  to predict. With an even smaller pool, the same requests could be
  preempted again and again, and throughput could fall below safe
  admission.
\end{itemize}

\end{document}
